% ============================================================
% Question Bank: ch08-03 Comparing Two Populations Visually
% Topic Title: Comparing Two Populations Visually
% CCSS: 7.SP.B.3
% Questions: 30 (18 MC + 7 SA + 5 Graphical)
% ============================================================

% --- Multiple Choice Questions (18) ---

\begin{practiceQuestion}{8-3-q01}{mc}
\begin{questionText}
Class A has a mean test score of $78$ and Class B has a mean test score of $85$. What can you conclude?
\end{questionText}
\choiceA{Every student in Class B scored higher than every student in Class A}
\choiceB{Class B scored higher on average than Class A}
\choiceC{Class A and Class B have the same scores}
\choiceD{Class A has more students than Class B}
\correctAnswer{B}
\explanation{The mean tells us the average performance. Class B scored higher on average, but individual scores may overlap.}
\end{practiceQuestion}

\begin{practiceQuestion}{8-3-q02}{mc}
\begin{questionText}
Two groups have overlapping dot plots. What does the overlap indicate?
\end{questionText}
\choiceA{The groups are identical}
\choiceB{Some data values from both groups are similar}
\choiceC{One group is always better than the other}
\choiceD{The data is incorrect}
\correctAnswer{B}
\explanation{Overlap means some members of each group have similar values, even if the centers differ.}
\end{practiceQuestion}

\begin{practiceQuestion}{8-3-q03}{mc}
\begin{questionText}
When comparing two populations visually, you should look at:
\end{questionText}
\choiceA{Only the center (mean or median)}
\choiceB{Only the spread (range or MAD)}
\choiceC{Both the center and the spread}
\choiceD{Only the largest and smallest values}
\correctAnswer{C}
\explanation{A complete comparison requires looking at both center (where the data clusters) and spread (how spread out it is).}
\end{practiceQuestion}

\begin{practiceQuestion}{8-3-q04}{mc}
\begin{questionText}
Group X has a median of $50$ and a range of $10$. Group Y has a median of $50$ and a range of $30$. Which statement is true?
\end{questionText}
\choiceA{Group X has more data values}
\choiceB{Group Y has a higher center}
\choiceC{Group Y is more spread out than Group X}
\choiceD{The groups are identical}
\correctAnswer{C}
\explanation{Both have the same median, but Group Y's larger range means its data is more spread out.}
\end{practiceQuestion}

\begin{practiceQuestion}{8-3-q05}{mc}
\begin{questionText}
Which display is best for comparing the overall shapes of two large data sets?
\end{questionText}
\choiceA{A single dot plot}
\choiceB{Two histograms}
\choiceC{A circle graph}
\choiceD{A frequency table with one column}
\correctAnswer{B}
\explanation{Histograms show the shape and distribution of large data sets, making it easy to compare two groups.}
\end{practiceQuestion}

\begin{practiceQuestion}{8-3-q06}{mc}
\begin{questionText}
Team A's scores: mean $72$, MAD $3$. Team B's scores: mean $74$, MAD $8$. Which team is more consistent?
\end{questionText}
\choiceA{Team A because its MAD is smaller}
\choiceB{Team B because its mean is higher}
\choiceC{Team A because its mean is lower}
\choiceD{Team B because its MAD is larger}
\correctAnswer{A}
\explanation{A smaller MAD means scores are closer to the mean, so Team A is more consistent.}
\end{practiceQuestion}

\begin{practiceQuestion}{8-3-q07}{mc}
\begin{questionText}
Two box plots are drawn on the same number line. Box plot A is entirely to the right of box plot B. What does this tell you?
\end{questionText}
\choiceA{Group A has more data}
\choiceB{All values in Group A are greater than all values in Group B}
\choiceC{Group A has no overlap with Group B}
\choiceD{Both B and C}
\correctAnswer{D}
\explanation{If box plot A is entirely to the right, every value in A exceeds every value in B, so there is no overlap.}
\end{practiceQuestion}

\begin{practiceQuestion}{8-3-q08}{mc}
\begin{questionText}
The dot plots for two classes overlap significantly. The difference in their means is $2$ points, and both have a MAD of $6$. What can you conclude?
\end{questionText}
\choiceA{The classes are very different}
\choiceB{The difference in means is meaningful}
\choiceC{The difference in means is small relative to the variability}
\choiceD{You cannot compare the classes}
\correctAnswer{C}
\explanation{The difference ($2$) is much smaller than the MAD ($6$), so the difference is not very meaningful — there is a lot of overlap.}
\end{practiceQuestion}

\begin{practiceQuestion}{8-3-q09}{mc}
\begin{questionText}
Which measure of variability is shown directly on a box plot?
\end{questionText}
\choiceA{Mean absolute deviation (MAD)}
\choiceB{Standard deviation}
\choiceC{Interquartile range (IQR)}
\choiceD{Mean}
\correctAnswer{C}
\explanation{A box plot shows the quartiles ($Q_1$ and $Q_3$), so you can read the IQR directly as $Q_3 - Q_1$.}
\end{practiceQuestion}

\begin{practiceQuestion}{8-3-q10}{mc}
\begin{questionText}
Group A: median $85$, IQR $8$. Group B: median $80$, IQR $20$. How do the groups compare?
\end{questionText}
\choiceA{Group A is higher and more consistent}
\choiceB{Group A is higher and less consistent}
\choiceC{Group B is higher and more consistent}
\choiceD{Both groups are equally consistent}
\correctAnswer{A}
\explanation{Group A has a higher median ($85$ vs.\ $80$) and a smaller IQR ($8$ vs.\ $20$), meaning it is higher and more tightly clustered.}
\end{practiceQuestion}

\begin{practiceQuestion}{8-3-q11}{mc}
\begin{questionText}
A dot plot for Group A is clustered tightly around $60$, while a dot plot for Group B is spread from $40$ to $80$. Both have a mean of $60$. Which statement is true?
\end{questionText}
\choiceA{Group B performed better}
\choiceB{Group A is more variable}
\choiceC{Group B is more variable}
\choiceD{The groups have different means}
\correctAnswer{C}
\explanation{Both means are $60$, but Group B's data spreads from $40$ to $80$, making it more variable.}
\end{practiceQuestion}

\begin{practiceQuestion}{8-3-q12}{mc}
\begin{questionText}
When is a dot plot a better choice than a box plot for comparing two groups?
\end{questionText}
\choiceA{When data sets are very large}
\choiceB{When you want to see every individual data point}
\choiceC{When you only want to compare medians}
\choiceD{When the data is categorical}
\correctAnswer{B}
\explanation{Dot plots show every data point, which is useful for small data sets where individual values matter.}
\end{practiceQuestion}

\begin{practiceQuestion}{8-3-q13}{mc}
\begin{questionText}
The difference between the means of two groups is $10$. The MAD for each group is $2$. How many MADs apart are the means?
\end{questionText}
\choiceA{$2$}
\choiceB{$5$}
\choiceC{$10$}
\choiceD{$20$}
\correctAnswer{B}
\explanation{$\frac{10}{2} = 5$ MADs apart.}
\end{practiceQuestion}

\begin{practiceQuestion}{8-3-q14}{mc}
\begin{questionText}
If the difference between two group means is more than $2$ times the MAD, the difference is generally considered:
\end{questionText}
\choiceA{Not meaningful}
\choiceB{Meaningless}
\choiceC{Meaningful}
\choiceD{Impossible to determine}
\correctAnswer{C}
\explanation{When the difference between means is more than $2$ MADs, there is little overlap and the difference is considered meaningful.}
\end{practiceQuestion}

\begin{practiceQuestion}{8-3-q15}{mc}
\begin{questionText}
Group X: mean $45$, range $12$. Group Y: mean $52$, range $12$. What can you say about these groups?
\end{questionText}
\choiceA{They have the same center and different spread}
\choiceB{They have different centers and the same spread}
\choiceC{They are identical}
\choiceD{Group X is more spread out}
\correctAnswer{B}
\explanation{The means differ ($45$ vs.\ $52$) but both have the same range ($12$), so they have different centers and similar spread.}
\end{practiceQuestion}

\begin{practiceQuestion}{8-3-q16}{mc}
\begin{questionText}
Two histograms show that Group A's data is skewed left and Group B's data is symmetric. Which measure of center is most appropriate to compare them?
\end{questionText}
\choiceA{Mean for both}
\choiceB{Median for both}
\choiceC{Mode for both}
\choiceD{Range for both}
\correctAnswer{B}
\explanation{The median is less affected by skew than the mean, so it is the best measure to compare a skewed and a symmetric distribution.}
\end{practiceQuestion}

\begin{practiceQuestion}{8-3-q17}{mc}
\begin{questionText}
Two groups have box plots that overlap completely. What does this suggest?
\end{questionText}
\choiceA{The groups have very similar distributions}
\choiceB{One group is clearly better than the other}
\choiceC{The data is unreliable}
\choiceD{The groups have no members in common}
\correctAnswer{A}
\explanation{Complete overlap in box plots means the groups have very similar medians, quartiles, and ranges.}
\end{practiceQuestion}

\begin{practiceQuestion}{8-3-q18}{mc}
\begin{questionText}
School A: median score $88$, MAD $4$. School B: median score $82$, MAD $4$. The difference in medians expressed in MADs is:
\end{questionText}
\choiceA{$1$ MAD}
\choiceB{$1.5$ MADs}
\choiceC{$2$ MADs}
\choiceD{$4$ MADs}
\correctAnswer{B}
\explanation{Difference in medians: $88 - 82 = 6$. In MADs: $\frac{6}{4} = 1.5$ MADs.}
\end{practiceQuestion}

% --- Short Answer Questions (7) ---

\begin{practiceQuestion}{8-3-q19}{sa}
\begin{questionText}
Group A: mean $65$, MAD $5$. Group B: mean $75$, MAD $5$. How many MADs apart are the means?
\end{questionText}
\correctAnswer{$2$ MADs}
\explanation{Difference: $75 - 65 = 10$. Divide by MAD: $\frac{10}{5} = 2$ MADs.}
\end{practiceQuestion}

\begin{practiceQuestion}{8-3-q20}{sa}
\begin{questionText}
Data set A: $10, 12, 14, 16, 18$. Find the range.
\end{questionText}
\correctAnswer{$8$}
\explanation{Range $= 18 - 10 = 8$.}
\end{practiceQuestion}

\begin{practiceQuestion}{8-3-q21}{sa}
\begin{questionText}
Data set: $4, 6, 7, 8, 10$. The mean is $7$. Find the MAD.
\end{questionText}
\correctAnswer{$1.6$}
\explanation{Deviations from mean: $|4-7|+|6-7|+|7-7|+|8-7|+|10-7| = 3+1+0+1+3 = 8$. MAD $= \frac{8}{5} = 1.6$.}
\end{practiceQuestion}

\begin{practiceQuestion}{8-3-q22}{sa}
\begin{questionText}
Team A: mean $60$, MAD $4$. Team B: mean $70$, MAD $3$. Is the difference in means meaningful? Explain briefly.
\end{questionText}
\correctAnswer{Yes; the means are about $2.5$ to $3.3$ MADs apart}
\explanation{Difference $= 10$. Using MAD of $4$: $\frac{10}{4} = 2.5$. Using MAD of $3$: $\frac{10}{3} \approx 3.3$. Both are greater than $2$, so the difference is meaningful.}
\end{practiceQuestion}

\begin{practiceQuestion}{8-3-q23}{sa}
\begin{questionText}
What does it mean when two box plots have a large gap between them on a number line?
\end{questionText}
\correctAnswer{The two groups have very different values with little or no overlap}
\explanation{A large gap indicates the groups' data ranges do not overlap much, suggesting a clear difference between the populations.}
\end{practiceQuestion}

\begin{practiceQuestion}{8-3-q24}{sa}
\begin{questionText}
Class A: median $90$, IQR $6$. Class B: median $83$, IQR $14$. Which class is more consistent?
\end{questionText}
\correctAnswer{Class A}
\explanation{Class A has a smaller IQR ($6$ vs.\ $14$), meaning the middle $50\%$ of scores are more tightly clustered.}
\end{practiceQuestion}

\begin{practiceQuestion}{8-3-q25}{sa}
\begin{questionText}
Two groups have means of $50$ and $54$, and both have MAD $= 8$. Express the difference in means as a multiple of MAD. Is the difference meaningful?
\end{questionText}
\correctAnswer{$0.5$ MADs; no, the difference is not meaningful}
\explanation{$\frac{54-50}{8} = 0.5$ MADs. Since $0.5 < 2$, the difference is small relative to the variability and not considered meaningful.}
\end{practiceQuestion}

% --- Graphical Questions (3 GMC + 2 GSA) ---

\begin{practiceQuestion}{8-3-q26}{gmc}
\begin{questionText}
The box plots below compare test scores for Class A and Class B.

\begin{center}
\begin{tikzpicture}[scale=0.1]
  \draw[thick,->] (49,0) -- (101,0);
  \foreach \x in {50,55,...,100} {
    \draw (\x,0.5) -- (\x,-0.5);
    \node[below,font=\small] at (\x,-0.5) {$\x$};
  }
  % Class A: min=60, Q1=70, med=78, Q3=85, max=95
  \draw[thick] (60,6) -- (95,6);
  \draw[thick,fill=funBlue!30] (70,4) rectangle (85,8);
  \draw[very thick,funBlue] (78,4) -- (78,8);
  \draw[thick] (60,4) -- (60,8);
  \draw[thick] (95,4) -- (95,8);
  \node[right,font=\small] at (96,6) {Class A};
  % Class B: min=50, Q1=58, med=65, Q3=72, max=80
  \draw[thick] (50,14) -- (80,14);
  \draw[thick,fill=funGreen!30] (58,12) rectangle (72,16);
  \draw[very thick,funGreen] (65,12) -- (65,16);
  \draw[thick] (50,12) -- (50,16);
  \draw[thick] (80,12) -- (80,16);
  \node[right,font=\small] at (81,14) {Class B};
\end{tikzpicture}
\end{center}

Which statement is best supported by the box plots?
\end{questionText}
\choiceA{Class B scored higher than Class A}
\choiceB{Class A has a higher median and a larger IQR}
\choiceC{Class A has a higher median and both classes have the same IQR}
\choiceD{Class A and Class B overlap, but Class A's median is higher}
\correctAnswer{D}
\explanation{Class A median $= 78$, Class B median $= 65$. Class A's range ($60$--$95$) overlaps with Class B's range ($50$--$80$), but Class A is higher overall.}
\end{practiceQuestion}

\begin{practiceQuestion}{8-3-q27}{gmc}
\begin{questionText}
The dot plots below show the number of push-ups completed by students in two gym classes.

\begin{center}
\begin{tikzpicture}[scale=0.55]
  % Class 1
  \node[left,font=\small\bfseries] at (4,2.5) {Class 1};
  \draw[thick,->] (4.5,0) -- (16.5,0);
  \foreach \x in {5,6,...,16} {
    \draw (\x,0.1) -- (\x,-0.1);
    \node[below,font=\tiny] at (\x,0) {$\x$};
  }
  % Class 1 data: 7(1),8(2),9(3),10(4),11(3),12(2),13(1)
  \fill[funBlue] (7,0.4) circle (0.15);
  \fill[funBlue] (8,0.4) circle (0.15); \fill[funBlue] (8,0.75) circle (0.15);
  \fill[funBlue] (9,0.4) circle (0.15); \fill[funBlue] (9,0.75) circle (0.15); \fill[funBlue] (9,1.1) circle (0.15);
  \fill[funBlue] (10,0.4) circle (0.15); \fill[funBlue] (10,0.75) circle (0.15); \fill[funBlue] (10,1.1) circle (0.15); \fill[funBlue] (10,1.45) circle (0.15);
  \fill[funBlue] (11,0.4) circle (0.15); \fill[funBlue] (11,0.75) circle (0.15); \fill[funBlue] (11,1.1) circle (0.15);
  \fill[funBlue] (12,0.4) circle (0.15); \fill[funBlue] (12,0.75) circle (0.15);
  \fill[funBlue] (13,0.4) circle (0.15);

  % Class 2
  \node[left,font=\small\bfseries] at (4,5.5) {Class 2};
  \draw[thick,->] (4.5,3.5) -- (16.5,3.5);
  \foreach \x in {5,6,...,16} {
    \draw (\x,3.6) -- (\x,3.4);
    \node[below,font=\tiny] at (\x,3.4) {};
  }
  % Class 2 data: 10(1),11(2),12(3),13(4),14(3),15(2),16(1)
  \fill[funGreen] (10,3.9) circle (0.15);
  \fill[funGreen] (11,3.9) circle (0.15); \fill[funGreen] (11,4.25) circle (0.15);
  \fill[funGreen] (12,3.9) circle (0.15); \fill[funGreen] (12,4.25) circle (0.15); \fill[funGreen] (12,4.6) circle (0.15);
  \fill[funGreen] (13,3.9) circle (0.15); \fill[funGreen] (13,4.25) circle (0.15); \fill[funGreen] (13,4.6) circle (0.15); \fill[funGreen] (13,4.95) circle (0.15);
  \fill[funGreen] (14,3.9) circle (0.15); \fill[funGreen] (14,4.25) circle (0.15); \fill[funGreen] (14,4.6) circle (0.15);
  \fill[funGreen] (15,3.9) circle (0.15); \fill[funGreen] (15,4.25) circle (0.15);
  \fill[funGreen] (16,3.9) circle (0.15);
\end{tikzpicture}
\end{center}

Which class has a higher center, and by approximately how many push-ups?
\end{questionText}
\choiceA{Class 1, by about $3$}
\choiceB{Class 2, by about $3$}
\choiceC{Class 2, by about $6$}
\choiceD{They have the same center}
\correctAnswer{B}
\explanation{Class 1 center $\approx 10$. Class 2 center $\approx 13$. Difference $\approx 3$ push-ups. Class 2 is higher.}
\end{practiceQuestion}

\begin{practiceQuestion}{8-3-q28}{gmc}
\begin{questionText}
The histograms show the distribution of weekly study hours for two groups.

\begin{center}
\begin{tikzpicture}[scale=0.55]
  % Group A
  \node[font=\small\bfseries] at (3.5,5) {Group A};
  \draw[thick,->] (0,0) -- (0,4.5);
  \draw[thick,->] (0,0) -- (7.5,0);
  \node[below,font=\small] at (3.5,-0.5) {Hours};
  \foreach \y in {1,2,3,4} { \draw (-0.1,\y) -- (0.1,\y); \node[left,font=\tiny] at (0,\y) {$\y$}; }
  \fill[funBlue!60] (0.5,0) rectangle (1.5,1);
  \fill[funBlue!60] (1.5,0) rectangle (2.5,2);
  \fill[funBlue!60] (2.5,0) rectangle (3.5,4);
  \fill[funBlue!60] (3.5,0) rectangle (4.5,3);
  \fill[funBlue!60] (4.5,0) rectangle (5.5,2);
  \fill[funBlue!60] (5.5,0) rectangle (6.5,1);
  \node[below,font=\tiny] at (1,0) {$2$};
  \node[below,font=\tiny] at (2,0) {$4$};
  \node[below,font=\tiny] at (3,0) {$6$};
  \node[below,font=\tiny] at (4,0) {$8$};
  \node[below,font=\tiny] at (5,0) {$10$};
  \node[below,font=\tiny] at (6,0) {$12$};
\end{tikzpicture}
\hspace{1cm}
\begin{tikzpicture}[scale=0.55]
  % Group B
  \node[font=\small\bfseries] at (3.5,5) {Group B};
  \draw[thick,->] (0,0) -- (0,4.5);
  \draw[thick,->] (0,0) -- (7.5,0);
  \node[below,font=\small] at (3.5,-0.5) {Hours};
  \foreach \y in {1,2,3,4} { \draw (-0.1,\y) -- (0.1,\y); \node[left,font=\tiny] at (0,\y) {$\y$}; }
  \fill[funGreen!60] (0.5,0) rectangle (1.5,3);
  \fill[funGreen!60] (1.5,0) rectangle (2.5,4);
  \fill[funGreen!60] (2.5,0) rectangle (3.5,2);
  \fill[funGreen!60] (3.5,0) rectangle (4.5,2);
  \fill[funGreen!60] (4.5,0) rectangle (5.5,1);
  \fill[funGreen!60] (5.5,0) rectangle (6.5,1);
  \node[below,font=\tiny] at (1,0) {$2$};
  \node[below,font=\tiny] at (2,0) {$4$};
  \node[below,font=\tiny] at (3,0) {$6$};
  \node[below,font=\tiny] at (4,0) {$8$};
  \node[below,font=\tiny] at (5,0) {$10$};
  \node[below,font=\tiny] at (6,0) {$12$};
\end{tikzpicture}
\end{center}

Which group tends to study more hours per week?
\end{questionText}
\choiceA{Group A, because most of its data is at higher values}
\choiceB{Group B, because most of its data is at higher values}
\choiceC{Group A, because it has more spread}
\choiceD{They study the same amount}
\correctAnswer{A}
\explanation{Group A's histogram peaks at $6$--$8$ hours, while Group B's peaks at $2$--$4$ hours. Group A tends to study more.}
\end{practiceQuestion}

\begin{practiceQuestion}{8-3-q29}{gsa}
\begin{questionText}
The box plots compare daily steps for two people over one month.

\begin{center}
\begin{tikzpicture}[scale=0.012]
  \draw[thick,->] (2900,0) -- (12100,0);
  \foreach \x in {3000,4000,...,12000} {
    \draw (\x,80) -- (\x,-80);
    \node[below,font=\small] at (\x,-80) {$\x$};
  }
  % Person A: min=4000, Q1=6000, med=7500, Q3=9000, max=11000
  \draw[thick] (4000,600) -- (11000,600);
  \draw[thick,fill=funBlue!30] (6000,400) rectangle (9000,800);
  \draw[very thick,funBlue] (7500,400) -- (7500,800);
  \draw[thick] (4000,400) -- (4000,800);
  \draw[thick] (11000,400) -- (11000,800);
  \node[right,font=\small] at (11200,600) {Person A};
  % Person B: min=5000, Q1=7000, med=8000, Q3=9500, max=10000
  \draw[thick] (5000,1400) -- (10000,1400);
  \draw[thick,fill=funGreen!30] (7000,1200) rectangle (9500,1600);
  \draw[very thick,funGreen] (8000,1200) -- (8000,1600);
  \draw[thick] (5000,1200) -- (5000,1600);
  \draw[thick] (10000,1200) -- (10000,1600);
  \node[right,font=\small] at (10200,1400) {Person B};
\end{tikzpicture}
\end{center}

Find the IQR for each person and state who is more consistent in their daily steps.
\end{questionText}
\correctAnswer{Person A: IQR $= 3{,}000$; Person B: IQR $= 2{,}500$. Person B is more consistent.}
\explanation{Person A: $Q_3 - Q_1 = 9{,}000 - 6{,}000 = 3{,}000$. Person B: $9{,}500 - 7{,}000 = 2{,}500$. Person B's smaller IQR means more consistent daily steps.}
\end{practiceQuestion}

\begin{practiceQuestion}{8-3-q30}{gsa}
\begin{questionText}
The dot plots show quiz scores for two classes (out of $10$).

\begin{center}
\begin{tikzpicture}[scale=0.6]
  % Class X
  \node[left,font=\small\bfseries] at (3.5,2) {Class X};
  \draw[thick,->] (3.5,0) -- (11.5,0);
  \foreach \x in {4,5,...,10} {
    \draw (\x,0.1) -- (\x,-0.1);
    \node[below,font=\small] at (\x,0) {$\x$};
  }
  % Class X data: 5(1),6(2),7(4),8(5),9(3),10(1) = 16 students
  \fill[funBlue] (5,0.35) circle (0.13);
  \fill[funBlue] (6,0.35) circle (0.13); \fill[funBlue] (6,0.65) circle (0.13);
  \fill[funBlue] (7,0.35) circle (0.13); \fill[funBlue] (7,0.65) circle (0.13); \fill[funBlue] (7,0.95) circle (0.13); \fill[funBlue] (7,1.25) circle (0.13);
  \fill[funBlue] (8,0.35) circle (0.13); \fill[funBlue] (8,0.65) circle (0.13); \fill[funBlue] (8,0.95) circle (0.13); \fill[funBlue] (8,1.25) circle (0.13); \fill[funBlue] (8,1.55) circle (0.13);
  \fill[funBlue] (9,0.35) circle (0.13); \fill[funBlue] (9,0.65) circle (0.13); \fill[funBlue] (9,0.95) circle (0.13);
  \fill[funBlue] (10,0.35) circle (0.13);

  % Class Y
  \node[left,font=\small\bfseries] at (3.5,5.5) {Class Y};
  \draw[thick,->] (3.5,3.5) -- (11.5,3.5);
  \foreach \x in {4,5,...,10} {
    \draw (\x,3.6) -- (\x,3.4);
  }
  % Class Y data: 4(2),5(3),6(4),7(3),8(2),9(1),10(1) = 16 students
  \fill[funGreen] (4,3.85) circle (0.13); \fill[funGreen] (4,4.15) circle (0.13);
  \fill[funGreen] (5,3.85) circle (0.13); \fill[funGreen] (5,4.15) circle (0.13); \fill[funGreen] (5,4.45) circle (0.13);
  \fill[funGreen] (6,3.85) circle (0.13); \fill[funGreen] (6,4.15) circle (0.13); \fill[funGreen] (6,4.45) circle (0.13); \fill[funGreen] (6,4.75) circle (0.13);
  \fill[funGreen] (7,3.85) circle (0.13); \fill[funGreen] (7,4.15) circle (0.13); \fill[funGreen] (7,4.45) circle (0.13);
  \fill[funGreen] (8,3.85) circle (0.13); \fill[funGreen] (8,4.15) circle (0.13);
  \fill[funGreen] (9,3.85) circle (0.13);
  \fill[funGreen] (10,3.85) circle (0.13);
\end{tikzpicture}
\end{center}

Find the mean score for each class and determine which class performed better on the quiz.
\end{questionText}
\correctAnswer{Class X mean $\approx 7.75$; Class Y mean $\approx 6.25$. Class X performed better.}
\explanation{Class X: $(5+12+28+40+27+10) \div 16 = 122 \div 16 \approx 7.63$. Class Y: $(8+15+24+21+16+9+10) \div 16 = 103 \div 16 \approx 6.44$. Class X has a higher mean.}
\end{practiceQuestion}
